\paragraph{\trMorita\trInvariance。}\label{sub:morita-invariance}
设 $A$、$B$ 是（\trLinear、\trGraded、\trUnital）\trCategory，$F_0,F_1\colon A\to B$
是（\trLinear、\trGraded、\trUnital）\trFunctor。设 $\tilde{A}\subset A$ 是一个子%
\trCategory，其包含\trFunctor{}是一个\trEquivalence。
\vkthmsync{6}
\begin{lemma}\label{thm:restriction-hochschild-iso}
\begingroup\emergencystretch=8em\hbadness=10000
限制映射给出\trHochschildCohomology{}群之间的一个\trIsomorphism{}：
$\mathit{HH}(A,B)\cong$ $\mathit{HH}(\tilde{A},B)$。
\par\endgroup
\end{lemma}
\begin{proof}
设 $\rho\colon\mathit{CC}(A,B)\to\mathit{CC}(\tilde{A},B)$ 是链层面上的限制。对任意
$X\in\Ob A$，选取 $\tilde{X}\in\Ob\tilde{A}$ 连同一个\trIsomorphism{}%
$u_X\colon X\to\tilde{X}$；对本已落在 $\tilde{A}$ 中的那些\trObject{} $X$，可以取
$\tilde{X}=X$，并令 $u_X$ 为\trIdentityMorphism。\textbf{Dennis 余迹}（Dennis
cotrace）是如下的\trChainMap：
\begin{equation*}
\begin{aligned}
&\chi\colon\mathit{CC}^{r+s}(\tilde{A},B)^s
  \longrightarrow\mathit{CC}^{r+s}(A,B)^s,\\
&(\chi T)^r(a_r,\ldots,a_1)
=G_1(u_{X_r})^{-1}\,T^d(u_{X_r}a_ru_{X_{r-1}}^{-1},\ldots,
  u_{X_1}a_1u_{X_0}^{-1})G_0(u_{X_0}).
\end{aligned}
\end{equation*}
这是一个\trHomotopyEquivalence，因为 $\rho\circ\chi=\id$，而
$\id-\chi\circ\rho=\partial\kappa+\kappa\partial$，其中
\begin{equation*}
\begin{split}
(\kappa T)(a_{r-1},\ldots,a_1)
&=\sum_n(-1)^{\clubsuit_{n+(r+s)}}\,G_1(u_{X_{r-1}})^{-1}\,
  T(u_{X_{r-1}}a_{r-1}u_{X_{r-1}}^{-1},\ldots,\\
&\hspace{14em}u_{X_{n+1}}a_{n+1}u_{X_n}^{-1},u_{X_n},a_n,\ldots,a_1).
\qedhere
\end{split}
\end{equation*}
\end{proof}

\begin{remark}\label{thm:morita-bimod-route}
或者，也可以走下面这条更抽象的路线。取\trGraded{} $A$-\trBimodule{}的阿贝尔%
\trCategory{}%
$\mathit{Bimod}^0(A)$；更精确地说，其\trObject{}是从 $(A^{\mathrm{opp}},A)$ 到%
\trGraded\trVectorSpace{}之\trCategory{}的\textbf{双\trFunctor}（bifunctor），而%
\trBimodule{}\trHomomorphism{}就是\trDegree{}为零的\trNaturalTransformation。任意一对%
\trFunctor{} $(G_0,G_1)$ 都给出一个\trBimodule，它由\trVectorSpace{}%
$\Hom_B(G_0(X_0),G_1(X_1))$ 组成。滥用记号起见，我们把这个\trBimodule{}径直记作 $B$，
而把察看\trFunctor{}对 $(\mathrm{Id}_A,\mathrm{Id}_A)$ 所得的相应\trObject{}记作
$A$。\trHochschild{}\trComplex{}可以解释成：它来自刚才提到的那个\trBimodule{}在
$\mathit{Bimod}^0(A)$ 中的一个\trProjective{}分解，这意味着
\begin{equation}\label{eq:hh-as-ext}
\mathit{HH}(A,B)\cong\operatorname{Ext}_{\mathit{Bimod}^0(A)}(A,B).
\end{equation}
这里我们考虑的是任意阶数的 $\operatorname{Ext}$ 群，而且还让 $B$ 带着其给定的%
\trGrading{}经受一切可能的\trSuspension，从而右边是双\trGraded{}的。出于初等的理由，
把双\trFunctor{}从 $A$ 限制到 $\tilde{A}$ 给出\trCategory{}之间的一个\trEquivalence{}%
$\mathit{Bimod}^0(A)\cong\mathit{Bimod}^0(\tilde{A})$，于是我们便得到
$\operatorname{Ext}$ 群之间的一个\trIsomorphism。这与代数的\trHochschild{}%
（上）\trHomology{}的经典\trMorita\trInvariance{}之间的关系（例如参见
[143, Theorem 9.5.6]）至此应当清楚了。事实上，我们的处境类似于由矩阵代数所提供的那一%
\trMorita\trInvariance{}的基本例子，而上文的初等论证即照搬 [96, Theorem 1.2.4]。
\end{remark}
